% modules/00-overview.tex — the manuscript's first module.
%
% Claim-bearing theorem environments here carry a \label, and every such label is a
% ledger node id in research/program/ledger.yaml with a matching `kind`. Structural
% labels (\section, \subsection, equations) are not nodes and need none.
%
% This file ships with orientation prose and no claims. A worked module holding an
% obstruction, an open question and a proved proposition is in example/modules/.
\documentclass[../main.tex]{subfiles}
\begin{document}

\section{Orientation}
\label{sec:overview}

Replace this section. State what this document studies, which question it is
organized around, and which conventions everything else rests on.

A convention this program fixes --- a sign, a scaling, a log base, which constant
absorbs what --- is stated as a \texttt{\textbackslash begin\{definition\}} under its
own \texttt{\textbackslash label} and carries a ledger node with
\texttt{kind: definition} and \texttt{status: defined}. Claims resting on it name it
in \texttt{depends\_on}, so changing it is a mathematical edit the checker can see.
There is no separate notation or normalization file: a second list would be a copy of
the dependency graph that nothing keeps honest.

An obstruction is a fence: a statement violating it is wrong by construction, so it is
read \emph{before} a proof is attempted or a numerical run is spent
(\texttt{CLAUDE.md} constraint 5). It is an ordinary ledger node ---
\texttt{kind: obstruction}, stated under a \texttt{\textbackslash label} in an
\texttt{obstruction} environment like any other claim --- and it lives here, rather
than in a registry of its own, because a fence is only useful next to the mathematics
it fences. State both halves: what it rules out, and what it leaves open.

\end{document}
